\documentclass[border=2pt,tikz]{standalone}
% Shared TikZ/pgfplots style for every figure in the manuscript.
\usepackage[T1]{fontenc}
\usepackage{mathpazo}
\usepackage{amsmath,amssymb}
\usepackage{tikz}
\usepackage{pgfplots}
\usepackage{pgfplotstable}
\pgfplotsset{compat=1.18}
\usetikzlibrary{arrows.meta,positioning,calc,backgrounds,fit,shapes.geometric}
\usepgfplotslibrary{fillbetween}

% --- palette: colour-blind safe, legible in greyscale -------------------
\definecolor{cHyd}{HTML}{1B6CA8}   % hydride / primary blue
\definecolor{cPro}{HTML}{C1492E}   % proton  / primary red
\definecolor{cAcc}{HTML}{7B5AA6}   % accent  / purple
\definecolor{cGold}{HTML}{B8860B}  % accent  / gold
\definecolor{cGrey}{HTML}{6E6E6E}
\definecolor{cInk}{HTML}{1A1A1A}
\definecolor{cBandR}{HTML}{C1492E}
\definecolor{cBandG}{HTML}{4F7942}

% --- base axis ----------------------------------------------------------
\pgfplotsset{
  icig/.style={
    width=68mm, height=52mm,
    scale only axis,
    axis line style={cInk, line width=0.45pt},
    tick style={cInk, line width=0.4pt},
    tick align=outside, tick pos=left,
    label style={font=\small, color=cInk},
    tick label style={font=\footnotesize, color=cInk},
    title style={font=\small\bfseries, color=cInk, align=left},
    legend cell align=left,
    legend style={draw=none, fill=none, font=\scriptsize,
                  inner sep=1pt, row sep=0.5pt},
    grid=both,
    major grid style={cGrey!22, line width=0.28pt},
    minor grid style={cGrey!12, line width=0.2pt},
  },
  % marker shorthands
  ptA/.style={mark=*, mark size=1.7pt, thick,
              mark options={fill=cHyd, draw=white, line width=0.4pt}},
  ptB/.style={mark=*, mark size=1.7pt, thick,
              mark options={fill=cPro, draw=white, line width=0.4pt}},
  hollowB/.style={mark=o, mark size=1.7pt, thick,
              mark options={draw=cPro, line width=0.7pt}},
}

% panel tag, e.g. \panel{a}
\newcommand{\panel}[1]{\textbf{(\textsf{#1})}}

\usetikzlibrary{arrows.meta,positioning,calc}
\begin{document}
\begin{tikzpicture}
% Overview figure, full text width.  Panel B is computed by
% analysis/export_fig_overview.py; panels A and C are schematics.
% Generated by analysis/export_fig_overview.py; do not edit.
\def\foFzero{-0.042086}
\def\foFtop{0.62}
\def\foYadhF{0.128726}
\def\foLightF{-0.261685}

\pgfplotsset{
  ov/.style={icig, scale only axis,
    label style={font=\fontsize{8.0}{9.6}\selectfont, color=cInk},
    tick label style={font=\fontsize{7.2}{8.6}\selectfont, color=cInk},
    title style={font=\fontsize{8.6}{10.4}\selectfont\bfseries, color=cInk,
                 align=left}},
}
\tikzset{
  ann/.style={font=\fontsize{7.2}{8.6}\selectfont, color=cInk, align=left,
              inner sep=0pt},
  annc/.style={font=\fontsize{7.2}{8.6}\selectfont, color=cInk, align=center,
              inner sep=0pt},
  hd/.style={font=\fontsize{8.6}{10.4}\selectfont\bfseries, color=cInk,
             anchor=north west, inner sep=0pt},
  box/.style={draw=#1, fill=#1!7, line width=0.45pt, rounded corners=1.5pt,
              inner sep=3pt, font=\fontsize{7.2}{8.8}\selectfont, text=cInk,
              align=left},
  st/.style={circle, draw=cInk, line width=0.4pt, inner sep=0pt,
             minimum size=5.0mm, font=\fontsize{6.4}{7.6}\selectfont},
  fl/.style={-{Latex[length=1.3mm,width=0.95mm]}, line width=0.42pt, cInk},
  lbl/.style={font=\fontsize{6.2}{7.2}\selectfont, color=cInk, inner sep=1.0pt},
  setarr/.style={line width=1.5pt, -{Latex[length=2.0mm,width=1.9mm]}},
}
\hyphenpenalty=10000 \exhyphenpenalty=10000

% =====================================================================
% (A) the experiment
% =====================================================================
\node[hd] at (0,0) {(\textit{\textbf{A}})\ \ two ratios, three unknowns};

\node[anchor=north west, inner sep=0pt, font=\fontsize{7.6}{9}\selectfont]
  (rx) at (0.05,-0.62)
  {$\mathrm{E}+\mathrm{S}_L\;
    \underset{k_{\mathrm{off}}}{\overset{k_{\mathrm{on}}}{\rightleftharpoons}}\;
    \mathrm{ES}_L\;\xrightarrow{\;k_L\;}\;\mathrm{E}+\mathrm{P}$};
\node[ann, anchor=north west, text width=46mm] (cmp)
  at ($(rx.south west)+(0,-0.16)$)
  {Substrates labeled with H, D and T at the transferred position ($L$)
   compete for one enzyme.};

\node[box=cHyd, anchor=north west, text width=42mm] (meas)
  at ($(cmp.south west)+(0,-0.20)$)
  {\textbf{measured}\quad $\KHT$, $\KDT$\\[1pt]
   ratios of $V/K$, for H and for D against T};
\node[box=cGrey, anchor=north west, text width=42mm] (unk)
  at ($(meas.south west)+(0,-0.16)$)
  {\textbf{unknown}\quad $x_{\mathrm H}$, $x_{\mathrm D}$ and $c$\\[1pt]
   intrinsic effects of the chemical step; reciprocal commitment
   $c=k_{\mathrm{off}}/k_{\mathrm T}$};

\node[anchor=north west, inner sep=0pt, font=\fontsize{7.6}{9}\selectfont]
  (map) at ($(unk.south west)+(0.02,-0.22)$)
  {$K_{L\mathrm T}=\dfrac{x_L\,(1+c)}{x_L+c}$};
\node[ann, anchor=west, text width=24mm] at ($(map.east)+(0.25,0)$)
  {two equations\\($L=\mathrm{H},\mathrm{D}$),\\three unknowns};

% =====================================================================
% (B) the compatible set as the range of a curve
% =====================================================================
\begin{axis}[ov, name=B, width=44mm, height=47mm,
  at={(6.25cm,-0.62cm)}, anchor=north west,
  xmin=0, xmax=1, ymin=-0.32, ymax=\foFtop,
  xlabel={masking fraction $(\KHT-1)/c$},
  ylabel={intrinsic offset $F^{\mathrm{int}}$},
  xtick={0,0.25,0.5,0.75,1}, xticklabels={0,0.25,0.5,0.75,1},
  ytick={-0.2,0,0.2,0.4,0.6},
  yticklabels={$-0.2$,$0$,$0.2$,$0.4$,$0.6$},
  clip=false,
  title={(\textit{\textbf{B}})\ \ one intrinsic pair per commitment}]
  % what the ground-channel model attains at these effect sizes
  \fill[cAcc!13] (axis cs:0,-0.32) rectangle (axis cs:1,\foFzero);
  \draw[cAcc, line width=0.8pt] (axis cs:0,\foFzero) -- (axis cs:1,\foFzero);
  \draw[cInk, line width=0.6pt, densely dashed] (axis cs:0,0) -- (axis cs:1,0);
  \node[ann, color=cAcc!85!black, anchor=north east] at (axis cs:0.985,-0.062)
    {gated model:\\$F<\Fz$};
  \node[ann, anchor=south east] at (axis cs:0.985,0.012) {$F=0$};

  \addplot[cPro, line width=1.1pt] table[x=m,y=F] {data/fo_mask_yadh.dat};
  \addplot[cHyd, line width=1.1pt, densely dashed] table[x=m,y=F]
    {data/fo_mask_light.dat};
  % the endpoints are approached as masking vanishes, not attained
  \addplot[only marks, mark=o, mark size=2.0pt, cPro, line width=0.7pt]
    table[x=m,y=F] {data/fo_obs_yadh.dat};
  \addplot[only marks, mark=o, mark size=2.0pt, cHyd, line width=0.7pt]
    table[x=m,y=F] {data/fo_obs_light.dat};
  \node[ann, color=cPro, anchor=north west] at (axis cs:0.30,0.215) {yeast ADH};
  \node[ann, color=cHyd, anchor=north west] at (axis cs:0.11,-0.255) {ecDHFR, $5\,^{\circ}$C};

  % the compatible sets, projected onto the offset axis
  \draw[setarr, cPro] (axis cs:1.075,\foYadhF) -- (axis cs:1.075,\foFtop);
  \draw[cPro, fill=white, line width=0.7pt] (axis cs:1.075,\foYadhF) circle (1.6pt);
  \draw[setarr, cHyd] (axis cs:1.16,\foLightF) -- (axis cs:1.16,\foFtop);
  \draw[cHyd, fill=white, line width=0.7pt] (axis cs:1.16,\foLightF) circle (1.6pt);
  \node[annc, rotate=90] at (axis cs:1.1175,0.40) {compatible sets};
\end{axis}

% =====================================================================
% (C) topology fixes the direction
% =====================================================================
\def\cx{12.25}
\node[hd] at (\cx-0.15,0) {(\textit{\textbf{C}})\ \ topology fixes the direction};

% a short offset axis with the set drawn on it
\newcommand{\faxis}[1]{%
  \draw[cGrey!60, line width=0.5pt] (#1) ++(0,-0.50) -- ++(0,1.00);
}
\node[annc, color=cGrey] at (\cx+2.75,-0.47) {set of $F$};

% ---- row 1: competition, no bypass -> concave, lower bound
\begin{scope}[shift={(\cx,-1.05)}]
  \node[st] (e) at (0,0) {E};  \node[st] (s) at (1.05,0) {ES};
  \node[lbl,anchor=west] (p) at (1.95,0) {P};
  \draw[fl] (s) -- node[lbl,above]{$k$} (p.west);
  \draw[fl] (s) -- node[lbl,above]{$k_{\mathrm{off}}$} (e);
  \faxis{2.75,0}
  \draw[setarr, cHyd] (2.75,-0.20) -- (2.75,0.50);
  \fill[cHyd] (2.75,-0.20) circle (1.5pt);
  \node[ann, anchor=north west, text width=46mm] at (-0.30,-0.40)
    {competing exit only:\\\textbf{lower bound}, no upper bound};
\end{scope}

% ---- row 2: bypass with fast dissociation -> convex, upper bound.  Without
% the k_off arrow, product would be certain after binding and K = 1 (blind):
% the convex map is the limit k_off >> k, k_b of the bypass map, not a
% separate finite network (no network in the enumeration has C = 0).
\begin{scope}[shift={(\cx,-2.55)}]
  \node[st] (e) at (0,0) {E};  \node[st] (s) at (1.05,0) {ES};
  \node[lbl,anchor=west] (p) at (1.95,0) {P};
  \draw[fl] (s) -- node[lbl,above]{$k$} (p.west);
  \draw[fl, line width=0.9pt] (s) -- node[lbl,above]{$k_{\mathrm{off}}$} (e);
  \draw[fl,cPro] (s) to[bend right=50] node[lbl,below,color=cPro]{$k_{b}$}
    (p.south west);
  \faxis{2.75,0}
  \draw[setarr, cPro] (2.75,0.20) -- (2.75,-0.50);
  \fill[cPro] (2.75,0.20) circle (1.5pt);
  \node[ann, anchor=north west, text width=46mm] at (-0.30,-0.62)
    {blind route, $k_{\mathrm{off}}\gg k,k_b$ (limit):\\\textbf{upper bound}, no lower bound};
\end{scope}

% ---- row 3: both -> no finite bound when both strengths are free
\begin{scope}[shift={(\cx,-4.20)}]
  \node[st] (e) at (0,0) {E};  \node[st] (s) at (1.05,0) {ES};
  \node[lbl,anchor=west] (p) at (1.95,0) {P};
  \draw[fl] (s) -- node[lbl,above]{$k$} (p.west);
  \draw[fl] (s) -- node[lbl,above]{$k_{\mathrm{off}}$} (e);
  \draw[fl,cPro] (s) to[bend right=50] node[lbl,below,color=cPro]{$k_{b}$}
    (p.south west);
  \faxis{2.75,0}
  \draw[line width=1.5pt, cAcc, {Latex[length=2.0mm,width=1.9mm]}-{Latex[length=2.0mm,width=1.9mm]}]
    (2.75,-0.50) -- (2.75,0.50);
  \node[ann, anchor=north west, text width=46mm] at (-0.30,-0.62)
    {both, strengths unknown: \textbf{no finite}\\\textbf{bound}; a bound on $\phi=k_b/k_{\mathrm T}$\\restores the lower bound};
\end{scope}

\end{tikzpicture}
\end{document}
